% !TEX program = lualatex
\documentclass[11pt,a4paper]{article}
\usepackage{fontspec}
\usepackage[russian,english]{babel}
\setmainfont{Tinos}
\setmonofont{DejaVu Sans Mono}
\usepackage[a4paper,margin=22mm]{geometry}
\usepackage{amsmath,amssymb,booktabs,longtable,tabularx,array,graphicx,hyperref,enumitem,xcolor}
\hypersetup{colorlinks=true,linkcolor=blue!50!black,urlcolor=blue!50!black}
\graphicspath{{figures/}}
\title{Экспериментальное приложение к работе\\«Фазово-секторная адресация цифровых разложений»}
\author{Иван Проскураков}
\date{Эксперименты 001--018}
\begin{document}
\selectlanguage{russian}
\maketitle

\section*{Назначение приложения}
Это приложение фиксирует ход исследования.  Его цель - показать не только итоговый положительный контур, но и отрицательные проверки, которые сузили пространство гипотез.

\begin{longtable}{p{0.07\linewidth}p{0.22\linewidth}p{0.30\linewidth}p{0.32\linewidth}}
\toprule
№ & Название & Проверяемая идея & Итог \\
\midrule
\endfirsthead
\toprule
№ & Название & Проверяемая идея & Итог \\
\midrule
\endhead
001 & Sector index and residue benchmark & Проверить, дают ли остатки по резонансным $K$ сжатие first-occurrence search. & Нет устойчивого lift. Секторный индекс работает, residue-only не работает. \\
002 & Sector tree, chord margins, drift features & Проверить quadtree-pruning, хордовые margin, drift/2D features. & Префиксное отсечение работает; margin полезен; drift features не дают pi-specific predictor. \\
003 & Sector transform matrices & Матрицы/пропорции между сетками. & Binary/quaternary exact; decimal/quaternary sparse; $1000$ нулей $\to L_4=1661$. \\
004 & Matrix target transform + resonance coordinates & Проверить совместимость target transform и $M=1221q+r$. & Contained precision $1.0$ для малых нулевых блоков; резонансный адрес прикрепляется к $M$. \\
005 & Finite matrix-sector locate MVP & Собрать первый конечный locate engine. & 19 запросов совпали с direct search; contained certificates строги. \\
006 & Resonant spiral memory & Проверить память спиральных рукавов. & Short-step controls дают lift, но резонансные масштабы near random. \\
007 & Walk-forward zero prediction & Прогнозировать следующие zero-block hits по прошлым. & Residue prediction не даёт robust lift относительно random/extreme baseline. \\
008 & Inverse tree & Обратить сертификат и изучить ветвление. & Обратное дерево содержит $10^M$ ветвей; нужна branch compression/external certificate. \\
009 & Quaternary resonant address read & Читать четверичные слова по резонансным указателям. & Указатели корректны, но простая выборка статистически near random. \\
010 & Closed finite locate contour & Проверить полный контур на всех паттернах длины 2--5. & $111100$ паттернов; exact match с direct search для всех найденных. \\
011 & Anchor-relative negative indexing & Проверить чтение с отрицательным offset вокруг якорей. & Работает как файловая адресация: anchor + occurrence + offset. \\
012 & 100-digit boundary tail & Проверить уникальность 100-значного хвоста. & В тестах хвост уникален; ожидаемые ложные совпадения пренебрежимо малы. \\
013 & Boundary continuation & Проверить 10 кандидатов на шаг + сертификат. & 100 будущих цифр восстановлены в backtest; ровно один кандидат принят. \\
014 & Algebraic suffix recovery & Извлечь суффикс из cross-base certificate. & При достаточной глубине сертификата кандидатный суффикс становится единственным. \\
015 & Next digit recovery & Проверить следующую цифру на границе $H=10^6$. & Получена цифра $3$; candidate count $1$. \\
016 & External certificate lemma validation & Массово проверить лемму внешнего сертификата. & Без сертификата $10^k$ кандидатов; с достаточно глубоким сектором - один суффикс. \\
017 & Next 10 digits by certificate & Восстановить следующий десятизначный блок на искусственной границе. & Для $H=2\cdot10^6$ получено $6121731251$, candidate count $1$. \\
018 & Machin external certificate & Использовать независимую формулу Мэчина как источник сертификата. & На искусственных границах внешний сертификат схлопнул $10^{10}$ кандидатов до одного блока. \\
\bottomrule
\end{longtable}

\section*{Главная экспериментальная граница}
Эксперименты 001--009 показывают, какие простые предикторы не дают результата.  Эксперименты 010--018 показывают положительный контур: конечный поиск и сертифицированное продолжение при наличии внешнего секторного сертификата.  Важное уточнение после Experiment 018: внешний сертификат является не мистическим внешним источником, а вычисляемым алгебраическим объектом.  В Experiment 018 он был получен из формулы Мэчина
\[
  \pi=16\arctan\frac15-4\arctan\frac1{239},
\]
а затем пропорционально переведён в десятичную сетку.  Поэтому итоговая граница такова:
\begin{itemize}
\item подтвержден finite locate внутри построенного горизонта;
\item подтверждена пропорциональная cross-base сертификация;
\item подтверждена лемма внешнего секторного сертификата;
\item не подтверждён автономный beyond-horizon predictor без внешнего сертификата.
\end{itemize}

\section*{Реплицируемость}
Архив содержит пакеты экспериментов 001--018, таблицы, графики и исходный модуль пропорциональной арифметики \texttt{pifs/proportion.py}.  Для удобства имена экспериментальных архивов имеют вид \texttt{pi\_fs\_experiment\_NNN\_package.zip}.

\end{document}
